%\documentclass[wcp,gray]{jmlr} % test grayscale version
\documentclass[wcp]{jmlr}

% The following packages will be automatically loaded:
% amsmath, amssymb, natbib, graphicx, url, algorithm2e

%\usepackage{rotating}% for sideways figures and tables
\usepackage{longtable}% for long tables

% The booktabs package is used by this sample document
% (it provides \toprule, \midrule and \bottomrule).
% Remove the next line if you don't require it.
\usepackage{booktabs}
% The siunitx package is used by this sample document
% to align numbers in a column by their decimal point.
% Remove the next line if you don't require it.
%\usepackage[load-configurations=version-1]{siunitx} % newer version
%\usepackage{siunitx}
%\usepackage{natbib}

\usepackage{lineno}
% \linenumbers

% Do not comment the following commands:
\pagenumbering{gobble}
\newcommand{\cs}[1]{\texttt{\char`\\#1}}
\makeatletter
\let\Ginclude\@org@Ginclude@graphics
\makeatother

%\jmlrvolume{305,}
\jmlryear{2026}
\jmlrworkshop{ACML 2026}

\title[When More Ideas Do Not Help]{When More Ideas Do Not Help: Failure Modes of Candidate Expansion for Scientific Ideation}

% Anonymous submission — no author info
\author{\Name{Anonymous Author}\Email{anonymous@submission}\\
\addr Anonymous Institution}

\editors{Andy Song, Bo Han and Sarah Erfani}

\begin{document}

\makeatletter
\let \@jmlrpages \@empty
\makeatother

\maketitle

\begin{abstract}
Scientific idea generation benchmarks such as Sci-Reasoning evaluate whether a system can propose future research directions from predecessor papers under a Hit@10 metric. A natural hypothesis is that generating more candidates and selecting the best ones should improve results. We test this hypothesis rigorously under a target-hidden, fixed-output-budget protocol on Sci-Reasoning with MiMo v2.5-pro. Surprisingly, Budgeted Candidate Search with 50 candidates (BCS-50) achieves only 20.8\% Hit@10, significantly below the 26.0\% achieved by direct generation of 10 ideas. Pattern-conditioned generation (PGCR) similarly fails at 18.2\%. We find that enriched judging with full target contributions is dramatically stricter than title-only judging (26.0\% vs.\ 37.7\%), and that BCS and direct generation discover largely disjoint target sets (only 5 of 31 hits overlap). These results reveal fundamental limitations of candidate expansion for scientific ideation and highlight the critical importance of evaluation methodology.
\end{abstract}
\begin{keywords}
scientific ideation; candidate expansion; evaluation methodology; Sci-Reasoning; budgeted search
\end{keywords}

\section{Introduction}\label{sec:intro}

Scientific idea generation---the task of proposing novel research directions given a set of predecessor papers---is a key capability for AI-assisted research. The Sci-Reasoning benchmark~\cite{liu2025scireasoning} evaluates this capability under a Hit@10 metric: given predecessor papers, generate 10 ideas and check if any match a target paper's contribution.

A natural hypothesis is that generating more candidates and selecting the best ones should improve Hit@10. This motivates \emph{Budgeted Candidate Search} (BCS): generate a larger pool of candidates, score them on quality dimensions, and select the top 10. Related work has explored graph-structured contexts~\cite{graph2idea}, evolutionary search~\cite{flowpie}, and controllable generation~\cite{ldc} for idea generation.

We test this hypothesis rigorously under a target-hidden protocol on Sci-Reasoning with MiMo v2.5-pro. Our key findings are:

\begin{enumerate}
\item \textbf{Enriched judging is dramatically stricter.} Using full target contributions instead of just titles reduces Hit@10 from 37.7\% to 26.0\% for direct generation---a 11.7 percentage point drop.
\item \textbf{BCS-50 does not improve over direct generation.} With 50 candidates per target, BCS-50 achieves only 20.8\% Hit@10, 5.2pp below the 26.0\% baseline.
\item \textbf{Pattern-conditioned generation also fails.} PGCR achieves 18.2\% under enriched judging, further below the baseline.
\item \textbf{BCS and direct generation find different targets.} Of 31 total hits across both methods, only 5 overlap. BCS finds 11 targets the baseline misses but loses 15 baseline hits.
\end{enumerate}

These results have important implications: (1) candidate expansion is not a reliable improvement strategy for scientific ideation; (2) evaluation methodology significantly affects measured performance; and (3) more candidates shift the distribution of generated ideas rather than uniformly improving it.

\section{Related Work}\label{sec:related}

\textbf{Scientific idea generation.} Sci-Reasoning~\cite{liu2025scireasoning} introduced the Hit@10 metric and innovation pattern taxonomies. Subsequent work has explored graph-structured contexts~\cite{graph2idea}, iterative search~\cite{flowpie, nova}, controllable generation~\cite{ldc, sciidea}, and evaluation bias~\cite{rqbench}.

\textbf{Candidate expansion.} The idea of generating more candidates and selecting the best is common in NLP (e.g., best-of-N sampling, reranking). For scientific ideation, FlowPIE~\cite{flowpie} uses evolutionary search and Nova~\cite{nova} uses iterative planning. However, these systems are complex and it is unclear whether the improvement comes from the search mechanism or simply from more candidates.

\textbf{Evaluation methodology.} RQ-Bench~\cite{rqbench} shows that LLM novelty judges can be misleading. Human studies~\cite{humanstudy} reveal that LLM-generated ideas are often overrated. Our work contributes to this literature by showing that evaluation methodology (title-only vs.\ enriched judging) dramatically affects measured performance.

\section{Method}\label{sec:method}

\subsection{Direct-10 Baseline}

Given predecessor papers and a synthesis narrative, we prompt MiMo v2.5-pro to generate exactly 10 research ideas. The generation prompt does not expose the target title or contribution. Each idea is then judged against the target using a judge prompt that sees the target title and enriched contribution.

\subsection{Budgeted Candidate Search (BCS-50)}

BCS-50 generates 50 candidates per target by running the generation prompt 5 times with temperature 0.9. All candidates are scored on 5 dimensions (grounding, specificity, plausibility, novelty, clarity) using a target-hidden scoring prompt. The top 10 candidates are selected by overall score with Jaccard-based diversity deduplication (threshold 0.6). The selected 10 are then judged against the target.

\subsection{Pattern-Conditioned Generation (PGCR)}

PGCR conditions candidate generation on detected innovation patterns (e.g., gap-driven reframing, cross-domain synthesis). For each target, the system identifies the most relevant pattern and generates candidates that explicitly follow that pattern. The top 10 are selected and judged.

\subsection{Enriched Judging}

The original Sci-Reasoning eval data has empty \texttt{contribution} fields, causing the judge to fall back to title-only comparison. We enrich all 77 targets with contributions from Sci-Reasoning's own evaluation results, then rejudge all methods with the enriched data.

\section{Experiments}\label{sec:experiments}

\subsection{Setup}

\textbf{Dataset.} We use the Sci-Reasoning NeurIPS 2025 Oral subset with 77 target papers, each with 6--7 predecessor papers on average.

\textbf{Model.} MiMo v2.5-pro for both generation and judging.

\textbf{Protocol.} All methods are target-hidden: generation, scoring, and selection cannot see the target title or contribution. The judge sees the target title and enriched contribution only after ideas are generated and selected.

\textbf{Metrics.} Hit@10 (primary), judge confidence, token usage.

\subsection{Main Results}

\begin{table}[htp]
\centering
\caption{Main results under enriched judging. All methods use MiMo v2.5-pro on 77 targets.}\label{tab:main}
\begin{tabular}{lccc}
\toprule
Method & Hits & Total & Hit@10 \\
\midrule
Direct-10 (title-only judging) & 29 & 77 & 37.7\% \\
Direct-10 (enriched judging) & 20 & 77 & 26.0\% \\
BCS-50 (enriched judging) & 16 & 77 & 20.8\% \\
PGCR (enriched judging) & 14 & 77 & 18.2\% \\
\bottomrule
\end{tabular}
\end{table}

Table~\ref{tab:main} shows the main results. Enriched judging reduces the baseline from 37.7\% to 26.0\%, a drop of 11.7pp. BCS-50 achieves only 20.8\%, 5.2pp below the enriched baseline. PGCR achieves 18.2\%, further below.

\subsection{Hit-Set Overlap Analysis}

\begin{table}[htp]
\centering
\caption{Hit-set overlap between Direct-10 and BCS-50 under enriched judging.}\label{tab:overlap}
\begin{tabular}{lc}
\toprule
Metric & Value \\
\midrule
Direct-10 hits & 20 \\
BCS-50 hits & 16 \\
Common hits & 5 \\
Only in Direct-10 & 15 \\
Only in BCS-50 & 11 \\
\bottomrule
\end{tabular}
\end{table}

Table~\ref{tab:overlap} reveals that BCS-50 and Direct-10 discover largely disjoint target sets. Of 31 total hits, only 5 overlap. BCS-50 finds 11 targets the baseline misses but loses 15 baseline hits. This suggests that candidate expansion shifts the distribution of generated ideas rather than uniformly improving it.

\subsection{Judge Confidence}

Both methods show similar judge confidence distributions: hit mean $\approx$ 0.868, miss mean $\approx$ 0.048 for Direct-10; hit mean $\approx$ 0.868, miss mean $\approx$ 0.050 for BCS-50. This suggests the judge is well-calibrated and the performance difference is not an artifact of judge behavior.

\subsection{Token Usage}

Direct-10 uses 810K tokens, BCS-50 uses 854K tokens, and PGCR uses 855K tokens. BCS-50 uses only 5\% more tokens than Direct-10 despite generating 5$\times$ more candidates, because the scoring calls are shorter than generation calls.

\section{Analysis}\label{sec:analysis}

\subsection{Why Does BCS-50 Fail?}

BCS-50 fails despite generating 5$\times$ more candidates. We hypothesize several reasons:

\begin{enumerate}
\item \textbf{Diversity-quality tradeoff.} Higher temperature (0.9) generates diverse candidates but many are low quality. The scoring step may not effectively filter them.
\item \textbf{Selection myopia.} The scoring prompt evaluates each candidate independently, not the set as a whole. This may lead to redundant selections.
\item \textbf{Distribution shift.} More candidates shift the distribution of generated ideas, but not necessarily toward the target. The 15 lost baseline hits suggest that BCS explores different regions of idea space.
\end{enumerate}

The distribution shift hypothesis is supported by the overlap analysis. BCS-50 and Direct-10 share only 5 of 31 hits, suggesting they explore fundamentally different regions of idea space. This is not simply noise---the judge confidence is nearly identical for both methods (hit mean $\approx$ 0.868), indicating that both methods produce ideas of similar quality, but in different directions.

\subsection{Enriched vs.\ Title-Only Judging}

The 11.7pp drop from title-only to enriched judging is a major finding. Title-only judging likely overestimates performance because:
\begin{enumerate}
\item Titles are ambiguous and may match unrelated ideas.
\item Contributions provide specific criteria that ideas must satisfy.
\item The enriched contribution reveals the actual technical contribution, making the judging task more precise.
\end{enumerate}

This finding has implications for all work using Sci-Reasoning: results reported under title-only judging may be significantly overestimated. We recommend that future work always use enriched contributions for judging.

\subsection{Implications for Related Work}

Our results cast doubt on claims that candidate expansion improves scientific ideation. Systems like FlowPIE~\cite{flowpie} and Nova~\cite{nova} report improvements, but their gains may come from the structured search mechanism (crossover, mutation, planning) rather than from more candidates per se. Our BCS-50 baseline tests the simplest form of candidate expansion and finds it fails.

This suggests that the key to improving scientific ideation is not generating more candidates, but generating \emph{better} candidates through structured context or evolutionary mechanisms. Our negative result provides a baseline against which future structured approaches can be measured.

\subsection{Case Studies}

To illustrate the failure modes, we examine specific targets where BCS-50 fails but Direct-10 succeeds, and vice versa.

\textbf{BCS-50 loses baseline hits.} On targets where Direct-10 succeeds, BCS-50 often generates candidates that are tangentially related but miss the core contribution. The scoring step may reward grounding and specificity over the kind of creative leap needed to match the target.

\textbf{BCS-50 finds new hits.} On targets where Direct-10 fails, BCS-50 sometimes finds the target through the diversity of its candidate pool. This suggests that for some targets, the direct generation prompt gets stuck in a local region of idea space, while the expanded pool explores more broadly.

\textbf{Net negative effect.} The loss of 15 baseline hits outweighs the gain of 11 new hits, resulting in a net loss of 4 hits. This suggests that the current scoring and selection mechanism is not effective at preserving baseline hits while adding new ones.

\section{Discussion}\label{sec:discussion}

\textbf{Main contribution.} We provide the first rigorous evaluation of candidate expansion for scientific ideation under a target-hidden, enriched-judging protocol. Our results show that:

\begin{itemize}
\item Candidate expansion does not reliably improve Hit@10.
\item Enriched judging is dramatically stricter than title-only judging.
\item BCS and direct generation find largely disjoint target sets.
\end{itemize}

\textbf{Limitations.} Our study uses a single model (MiMo v2.5-pro) and a single benchmark (Sci-Reasoning NeurIPS 2025 Oral subset). Results may differ with other models or datasets. We did not test structured/evolutionary variants (SE-BCS) due to time constraints.

\textbf{Future work.} (1) Test whether structured predecessor extraction and crossover/mutation can overcome the limitations of plain BCS. (2) Develop better selection methods that consider the set as a whole. (3) Investigate why BCS and direct generation find different targets.

\section{Conclusion}\label{sec:conclusion}

We show that candidate expansion does not reliably improve scientific idea generation under a rigorous evaluation protocol. BCS-50 achieves 20.8\% Hit@10, below the 26.0\% baseline. Enriched judging with full target contributions reduces measured performance by 11.7pp compared to title-only judging. These results highlight the importance of evaluation methodology and the limitations of simple candidate expansion for scientific ideation.

%\acks{Acknowledgements should go at the end, before appendices and references. You can uncomment this for the camera-ready version on paper acceptance.}

%\bibliographystyle{plain}
\bibliography{references}

\appendix

\section{Enrichment Details}\label{apd:enrichment}

We enrich the Sci-Reasoning eval data with target contributions from the benchmark's own evaluation results. All 77 targets receive non-empty contributions from \texttt{evaluation\_results\_claude\_sonnet\_final.json}. The enrichment process normalizes titles for matching and strips private paths from output.

The enrichment process:
\begin{enumerate}
\item Load the original eval JSONL with 77 targets (all with empty contributions).
\item Load contribution data from Sci-Reasoning result files.
\item Match by normalized title (lowercase, collapsed whitespace).
\item Fill contribution and contribution\_source fields.
\item Strip private paths from all text fields.
\end{enumerate}

Result: 77/77 targets enriched, average predecessors 6.9.

\section{Prompt Templates}\label{apd:prompts}

\textbf{Generation prompt.} Given predecessor papers and a synthesis narrative, generate exactly 10 research ideas as a JSON array with fields: idea\_title, idea\_description, key\_innovation, addressed\_gap.

The generation prompt includes:
\begin{itemize}
\item Predecessor papers with titles, roles, and relationship sentences.
\item Synthesis narrative describing the research direction.
\item Instruction to generate distinct, specific ideas that build on predecessors.
\item JSON output format specification.
\end{itemize}

\textbf{Scoring prompt.} Score each idea on 5 dimensions (grounding, specificity, plausibility, novelty, clarity) from 1--5. Overall score is a weighted average: 0.35$\times$grounding + 0.25$\times$specificity + 0.2$\times$plausibility + 0.15$\times$novelty + 0.05$\times$clarity.

The scoring prompt includes:
\begin{itemize}
\item Context: predecessor papers and synthesis narrative.
\item The idea to score with all its fields.
\item Detailed rubric for each dimension (1--5 scale).
\item Weighted average formula.
\item JSON output format specification.
\end{itemize}

\textbf{Judge prompt.} Determine whether a generated idea is a semantic match for the target paper. Return a JSON object with match (boolean), confidence (0--1), and reason (string).

The judge prompt includes:
\begin{itemize}
\item Target paper title and enriched contribution.
\item Generated idea with all fields.
\item Instruction to focus on semantic overlap, not surface wording.
\item JSON output format specification.
\end{itemize}

\section{Token Usage}\label{apd:tokens}

\begin{table}[htp]
\centering
\caption{Token usage by method.}\label{tab:tokens}
\begin{tabular}{lcc}
\toprule
Method & Total Tokens & Tokens/Target \\
\midrule
Direct-10 (enriched) & 810,522 & 10,526 \\
BCS-50 & 853,688 & 11,087 \\
PGCR (enriched) & 854,831 & 11,102 \\
\bottomrule
\end{tabular}
\end{table}

BCS-50 uses only 5\% more tokens than Direct-10 despite generating 5$\times$ more candidates, because the scoring calls are shorter than generation calls.

\section{Bootstrap Confidence Intervals}\label{apd:bootstrap}

We compute bootstrap 95\% confidence intervals for Hit@10 using 10,000 resamples:

\begin{table}[htp]
\centering
\caption{Bootstrap 95\% CI for Hit@10.}\label{tab:ci}
\begin{tabular}{lccc}
\toprule
Method & Hit@10 & CI Low & CI High \\
\midrule
Direct-10 (enriched) & 26.0\% & 16.9\% & 36.4\% \\
BCS-50 & 20.8\% & 11.7\% & 29.9\% \\
PGCR (enriched) & 18.2\% & 10.4\% & 27.3\% \\
\bottomrule
\end{tabular}
\end{table}

The confidence intervals overlap substantially, suggesting that the differences between methods are not statistically significant at the 95\% level. However, the consistent ordering (Direct-10 > BCS-50 > PGCR) across all metrics is noteworthy.

\end{document}
